% Generated draft: requires mathematical and source review.
\documentclass[11pt]{article}
\usepackage[margin=25mm]{geometry}
\usepackage{fontspec}
\setmainfont{Latin Modern Roman}
\usepackage{amsmath,amssymb,mathtools,unicode-math}
\setmathfont{Latin Modern Math}
\usepackage{enumitem,needspace,xurl,hyperref}
\setlength{\parindent}{0pt}
\setlength{\parskip}{5pt}
\setlength{\emergencystretch}{3em}
\newcommand{\N}{\mathbb{N}}
\newcommand{\Z}{\mathbb{Z}}
\newcommand{\Q}{\mathbb{Q}}
\newcommand{\R}{\mathbb{R}}
\newcommand{\C}{\mathbb{C}}
\newcommand{\F}{\mathbb{F}}
\newcommand{\A}{\mathbb{A}}
\newcommand{\PP}{\mathbb{P}}
\newcommand{\E}{\mathbb{E}}
\newcommand{\eps}{\varepsilon}
\newcommand{\dd}{\,\mathrm{d}}
\begin{document}
\section*{UnsolvedMath ID: 8800093}
\textbf{Source problem code:} catalogue entry
\subsection{Definitions and mathematical statement}
The elliptic-curve discrete-logarithm problem takes points $P,Q$ on $E/\mathbb F_q$, with $P$ of known finite order $r$ and $Q\in\langle P\rangle$, and asks for $a\in\mathbb Z/r\mathbb Z$ such that $Q=aP$.

The source asks whether its lifting attack can be improved by permitting the degree $d=[K:\mathbb Q]$ of an auxiliary number field to grow. A typical lifting framework chooses a number field with a prime whose residue field supports the finite-field curve and lifts relevant geometric/arithmetic data. The precise Huang--Raskind-style construction and the role of Shafarevich--Tate arithmetic are not provided, so no specific lift existence or attack formula is asserted.

Determine whether a nonconstant choice of $d$ as the problem parameters grow gives a genuine improvement over the intended fixed-degree method. A comparison must specify the baseline complexity and success probability and include the cost of representing and constructing $K$, its arithmetic, discriminant/height growth, and every remaining attack phase. Allowing uncounted higher-degree arithmetic is not an algorithmic improvement. The source does not prescribe a degree-growth regime, computational model, heuristic/conditional assumptions, or even the exact baseline lifting algorithm, so no numerical speedup threshold can be formalized from the record alone.
\subsection{Short English statement}
Can growing the auxiliary number-field degree improve the source's elliptic-curve discrete-log lifting attack after all field-construction and arithmetic costs are counted? The attack details and comparison baseline are missing.
\subsection{Sources}
\begin{itemize}
\item Ming-Deh Huang, IPAM \emph{Workshop on Number Theory and Cryptography --- Open Problems} (2006), attachment 6673, PDF page 9, schedule source \url{https://www.ipam.ucla.edu/programs/workshops/workshop-i-number-theory-and-cryptography-open-problems/?tab=schedule}; accessed 2026-07-29. The supplied provenance records a PDF/legacy-PPT visual audit.
\item Huang--Raskind-style lifting and Shafarevich--Tate methods are named in the supplied summary without a full publication citation or exact attack statement.
\item Local registry \texttt{data/unsolvedmath/problems.json}, permanent UnsolvedMath ID $8800093$, source code catalogue entry, \emph{Growing-degree improvements to the lifting attack}. Dataset attribution: Ulam AI and the UnsolvedMath Contributors.
\end{itemize}
\end{document}
